\section{Экспериментальное исследование}
\subsection{Визуальный анализ и характеристика данных}

При исследовании хаотических систем критически важно различать два типа различий, 
возникающих при варьировании управляющих параметров. Первый — \textit{Топологические различия}
 в точках бифуркации, где происходит изменение топологии динамической системы.
  (например, переход от хаотического режима к устойчивой периодической орбите).
   Второй — \textit{Пространственные  различия}, наблюдаемые в интервалах между бифуркациями семейста рядов Лоренца,
    когда аттрактор сохраняет топологическую структуру, но меняются его метрические свойства.

На рис.~\ref{fig:lorenz_comparison} представлены фрагменты траекторий, иллюстрирующие 
именно пространственные различия. Поскольку все параметры \(r\) выбраны в зоне
 устойчивого хаоса (\(24.74 < r < 50\)), топологическая форма аттрактора остаётся инвариантной.

\begin{figure}[htbp]
    \centering
    \includegraphics[width=\textwidth]{rows.png} % Замените путь!
    \caption{Сравнение участков траекторий системы Лоренца при различных значениях параметра \(r\).
     Все ряды принадлежат одной области между бифуркациями, поэтому их топологическая структура сохраняется. 
     Визуально заметно: чем больше отклонение \(r\) от эталонного значения 28, тем интенсивнее расходятся
     траектории.}
    \label{fig:lorenz_comparison}
\end{figure}

Для оценки сходства были вычислены стандартные метрики: динамическое время-образование
 (DTW) и корреляция Пирсона (табл.~\ref{tab:lorenz_metrics}).

\begin{table}[htbp]
\centering
\caption{Метрики сходства временных рядов системы Лоренца для различных значений параметра \(r\) (эталон: \(r=28\)).}
\label{tab:lorenz_metrics}
\begin{tabular}{|r|r|r|r|r|}
\hline
\(r\) & DTW & Евклидово расстояние & Корреляция Пирсона & Корреляция Спирмена \\
\hline
28.0001 & 16386.42 & 959.55 & 0.2574 & 0.2819 \\
\hline
27.9900 & 16994.62 & 1045.52 & 0.1277 & 0.1439 \\
\hline
28.0100 & 18952.64 & 1042.86 & 0.1094 & 0.1254 \\
\hline
28.1000 & 17823.66 & 1009.60 & 0.1885 & 0.2159 \\
\hline
27.9000 & 19479.54 & 1101.37 & -0.0099 & 0.0153 \\
\hline
30.0000 & 20504.81 & 1095.14 & 0.0933 & 0.1106 \\
\hline
31.0000 & 20146.69 & 1079.90 & 0.1104 & 0.1290 \\
\hline
32.0000 & 18226.30 & 1143.92 & 0.0129 & 0.0289 \\
\hline
33.0000 & 18633.11 & 1095.91 & 0.1271 & 0.1406 \\
\hline
34.0000 & 20951.99 & 1157.70 & 0.0485 & 0.0536 \\
\hline
35.0000 & 22030.53 & 1157.61 & 0.0622 & 0.0731 \\
\hline
36.0000 & 25929.73 & 1222.55 & -0.0184 & -0.0067 \\
\hline
37.0000 & 24355.96 & 1236.83 & -0.0242 & -0.0232 \\
\hline
38.0000 & 27064.00 & 1208.58 & 0.0638 & 0.0723 \\
\hline
39.0000 & 26136.93 & 1173.34 & 0.1090 & 0.1124 \\
\hline
\end{tabular}
\end{table}

На первый взгляд, визуальный (рис.~\ref{fig:lorenz_comparison}) и количественный 
(табл.~\ref{tab:lorenz_metrics}) анализ даёт чёткий ответ: ряд с \(r=28.0001\) практически
совпадает с эталоном на протяжении ~2000 шагов и демонстрирует минимальное DTW-расстояние
 с максимальной корреляцией, что делает его интуитивно предпочтительным кандидатом для аугментации.

Однако эта оценка, основанная на глобальных метриках, не подтверждается экспериментально.
 Дальнейший анализ показывает, что ряд при \(r=28.0001\), несмотря на кажущееся превосходство, 
 оказывается малоэффективным для расширения выборки(по крайней мере,на небольших размерах выборки).
  В то же время ряд с \(r=27.99\), формально 
 менее схожий по стандартным метрикам, демонстрирует высокую прогностическую ценность,
  в некоторых случаях превосходя эффект 
 от аугментации данными самого целевого процесса.

Это противоречие указывает на ключевое ограничение: \textit{метрики DTW и корреляция не
 способны оценить пригодность ряда для расширения обучающей выборки}. 
Это связано с тем,что эти метрики не предназначены для сравнения 
близких нелинейных систем из одного параметрического семейства.
Таким образом, данные табл.~\ref{tab:lorenz_metrics} и рис.~\ref{fig:lorenz_comparison} 
следует рассматривать как подтверждение того,что задача аугментации в предлагаемом методе 
требует разработки или апробации новой метрики.

\subsection{Влияние аугментации на качество прогноза}
\subsubsection{Описание эксперимента}
Основной целью данного эксперимента являлась оценка влияния размера и
состава обучающей выборки на качество многошагового прогнозирования. Экспериментальная
 установка была спроектирована для систематического исследования эффекта аугментации данных
  как из самого целевого процесса, так и из внешних, но схожих по динамике, систем.

В качестве целевого временного ряда $Y^{(0)}$ использовалась первая координата
 системы Лоренца, сгенерированная при канонических параметрах ($\sigma = 10$, $r = 28$, $b = 8/3$).
 Исходная обучающая выборка $Y_1^{(0)}$ состояла из 10 000 наблюдений. Для каждого из исследуемых
  вспомогательных рядов $Y^{(m)}$, сгенерированных с параметром $r$, незначительно отличающимся от
   базового, проводилась следующая итеративная процедура: к исходной выборке $Y_1^{(0)}$
   последовательно добавлялись сегменты данных из $Y^{(m)}$ до тех пор, пока общий размер
   объединенной выборки $\mathcal{D}_{\text{train}}$ не достигал 40 000 точек. Контрольный
   случай представлял собой аналогичное увеличение объема данных, но с использованием наблюдений
    из самого целевого ряда (т.е. $Y_1^{(0)}$ расширялся до 40 000 точек без привлечения внешних данных).

На каждом шаге увеличения размера выборки для полученного обучающего
 множества $\mathcal{D}_{\text{train}}$ выполнялось прогнозирование
 на тестовом наборе на горизонт в 10 шагов вперед. Для оценки качества прогноза 
 вычислялись три ключевые метрики: \textbf{доля непрогнозируемых точек (NP)},
  характеризующая уверенность модели; \textbf{среднеквадратичная ошибка (RMSE)}
   и \textbf{средняя абсолютная процентная ошибка (MAPE)}, оцениваемые на прогнозируемых точках.

Идентификация непрогнозируемых точек производилась с помощью практическог
о метода, основанного на кластеризации множества $\hat{S}_{t+h}$ алгоритмом DBSCAN
и проверке условия доминирования крупнейшего кластера.

\subsubsection{Анализ результатов}
\begin{figure}[htbp]
    \centering
    \includegraphics[width=0.8\textwidth]{metrics_vs_general_size_best_experiments.png} % Замените путь!
    \caption{Зависимость метрик от суммарного размера обучающей выборки для
    случаев эффективного расширения (r = 28.1, 28.0001, -27.9, -27.99,28).}
    \label{fig:main_exp}
\end{figure}

\paragraph{Анализ случаев улучшения метрик (Рисунок \ref{fig:main_exp}):} Для вспомогательных
 рядов с параметрами `r`, близкими к 28,(r=27.99,r=27.9) наблюдается значительное (до 6 раз) снижение доли
 непрогнозируемых точек, а также снижение RMSE в два раза. При этом метрики точности RMSE
 и MAPE остаются на уровне, сопоставимом с базовым случаем, или улучшаются. Это подтверждает
 гипотезу о том, что добавление данных из близких динамических систем повышает робастность
  модели без существенной потери точности.При этом для значения r = 28.0001
  количество непрогнозируемых точек падает логарифмически,но метрики
   качества на большинстве измерений ухудшаются.


\begin{figure}[htbp]
    \centering
    \includegraphics[width=0.8\textwidth]{unlucky_experiments.png} % Замените путь!
    \caption{Зависимость метрик от суммарного размера обучающей выборки для
    случаев неэффективного расширения (r = 28.01, 31, 38, 34,31,30). }
    \label{fig:unlucky_exp}
\end{figure}

\paragraph{Анализ случаев ухудшения метрик (Рисунок \ref{fig:unlucky_exp}):}
 Когда параметр `r` вспомогательного ряда значительно отличался от базового,
  снижение количества непрогнозируемых точек происходило ценой резкого роста
   ошибок RMSE и MAPE. Это указывает на деградацию качества прогноза при
   смешивании данных из систем с существенно различной динамикой. Однако,
    не все случаи с близкими параметрами приводят к улучшению:
    например, при добавлении данных с r=28.01 предсказание не улучшается
     и остается на уровне базовой выборки из 10000 точек, без снижения MAPE или RMSE.
Также важен тот факт,что вне зависимости от выбора вспомогательного ряда,
  увеличение объема обучающей выборки всегда приводит к снижению доли
   непрогнозируемых точек при устремлении к бесконечности размера добавленного ряда.Это видно на примере 
   r = 34,r=31,r=30-со увеличением размера этих рядов в обучающей выборке
    число непрогнозируемых точек начинает снижаться.
     Это связано с тем,что со временем z-вектора добавленного ряда начинают 
   доминировать над z-векторами оригинального ряда.Но если  ряд близок по параметру,
    то новые z-вектора будут согласоваться
    с оригинальными.Это подтверждается графиками \ref{fig:main_exp} и \ref{fig:unlucky_exp}: 
    даже при добавлении небольшого объема данных (5000 точек) к основной выборке (10000 точек) 
    количество непрогнозируемых точек сразу начинает снижаться.

\subsection{Исследование зависимости эффекта аугментации от расхождения параметров.}
\subsubsection{Описание эксперимента}
Экспериментальная установка была определена следующим образом. В
 качестве целевого ряда $Y^{(0)}$ использовалась траектория системы Лоренца с каноническим параметром 
 $r_0=28$. Из нее была сформирована базовая обучающая выборка $\mathcal{D}_{\text{base}}$, состоящая
  из первых $N_{\text{base}} = 5000$ наблюдений. Параллельно был сгенерирован набор вспомогательных временных
   рядов $\{Y^{(r)}\}$ для каждого значения параметра $r$ из исследуемого множества $\mathcal{R}$.
Ключевым шагом эксперимента стало формирование гибридных обучающих выборок $\mathcal{D}_{\text{train}}^{(r)}$ 
для каждого $r \in \mathcal{R}$. Каждая такая выборка создавалась путем объединения базовой выборки $\mathcal{D}_{\text{base}}$
 с сегментом из первых $N_{\text{aux}} = 5000$ наблюдений соответствующего вспомогательного ряда $Y^{(r)}$:
\[
\mathcal{D}_{\text{train}}^{(r)} = \mathcal{D}_{\text{base}} \cup \{y_0^{(r)}, y_1^{(r)}, \dots, y_{N_{\text{aux}}-1}^{(r)}\}.
\]
Таким образом, для всех $r$ размер гибридной выборки оставался постоянным и равным 10000 точкам.
 Наконец, для каждой выборки $\mathcal{D}_{\text{train}}^{(r)}$ производилось прогнозирование на
  тестовом наборе (из ряда $Y^{(0)}$) и вычислялись значения функционалов качества: $RMSE(r)$, $MAPE(r)$ и $NP(r)$.


Результаты, представленные как функции от параметра \(r\), сравнивались с двумя контрольными уровнями:
\begin{enumerate}
    \item Качество прогноза на малой выборке, \(\text{RMSE}(\mathcal{D}_{\text{base}})\) и
     \(\text{NP}(\mathcal{D}_{\text{base}})\).
    \item Качество прогноза на «чистой» выборке эквивалентного размера, состоящей
     только из данных целевого ряда, \(\mathcal{D}_{\text{target\_large}} = \{y_0^{(0)}, \dots, y_{9999}^{(0)}\}\).
\end{enumerate}
\begin{figure}[htbp]
    \centering
    \includegraphics[width=0.8\textwidth]{Different_r_comparison_with_points.png} % Замените путь!
    \caption{Сравнение метрик при добавлении 5000 точек из ряда с различным
    параметром r к 5000 точкам основного ряда.}
    \label{fig:different_r}
\end{figure}
\subsubsection{Анализ эксперимента}

Результаты (Рисунок \ref{fig:different_r}) показывают, что существует узкая окрестность (26,30)
 вокруг `r=28`, в которой добавление вспомогательных рядов приводит к
  улучшению метрик по сравнению с базовой выборкой из 5000 точек. Эта окрестность
  никак не связана с бифуркационной картиной.
  В лучших случаях достигается снижение RMSE и NP в 1.5 раза.
   Однако даже в лучших случаях качество прогноза уступает результатам,
    полученным на выборке из 10000 точек исходного ряда.


\subsection{Разработка и валидация критерия отбора рядов-доноров}

\subsubsection{Описание эксперимента}

В данном эксперименте была проведена апробация гибридной метрики: первый этап отбора
 использует сравнение статистических инвариантов динамики (оператор Перрона-Фробениуса),
  а второй этап отбора —оценивает геометрическое расстояние между аттракторами. Такой подход позволяет
   фильтровать не только ряды с различной эргодической структурой, но и устранять выбросы 
   с формально похожими операторами, но непригодную для метода мотивов геометрией аттрактора.

\paragraph{Алгоритм вычисления метрик схожести.}

\textbf{Первый этап отбора:оператор Перрона-Фробениуса.}
Теоретической основой для сравнения динамических систем служит оператор Перрона-Фробениуса,
 который описывает эволюцию плотностей вероятности состояний системы во времени. 
 Поскольку аналитическое выражение для этого оператора в сложных системах недоступно, 
 для его оценки применяется численный подход, заключающийся в построении эмпирической 
 конечномерной матрицы вероятностей перехода.

Алгоритм ее построения состоит из следующих шагов:

\begin{enumerate}
    \item \textbf{Дискретизация фазового пространства.} Для корректного сравнения 
    динамик оба временных ряда, эталонный $Y^{(1)}$ и ряд-кандидат $Y^{(2)}$, приводятся
     к дискретному представлению на общей сетке. Определяется глобальный диапазон амплитуд 
     $[V_{\min}, V_{\max}]$, который затем разбивается на $K$ непересекающихся интервалов 
     (состояний) $\{C_1, \dots, C_K\}$.
     Число состояний $K$ было выбрано из условия обеспечения статистической надежности
       матрицы перехода. Исходя истранномуз эмпирической оценки иметь в среднем не менее 10 наблюдений
        на каждый из $K^2$ типов переходов для выборки $N=10000$, получаем
         $K \le \sqrt{10000/10} \approx 31.6$. Таким образом, было установлено $K=32$.
 Каждое
      непрерывное значение в рядах заменяется на номер состояния, что приводит к 
      формированию двух последовательностей дискретных состояний, 
      $\mathcal{S}^{(1)}$ и $\mathcal{S}^{(2)}$.

    \item \textbf{Построение эмпирической матрицы перехода.} Для каждой 
    последовательности $\mathcal{S}$ строится стохастическая матрица $\mathbf{T}$ 
    размером $K \times K$, где элемент $T_{ij}$ оценивает условную вероятность перехода
    из состояния $C_i$ в $C_j$ за один шаг. Эта матрица является конечномерной аппроксимацией
     оператора Перрона-Фробениуса.

    \item \textbf{Вычисление расхождения.} Схожесть систем оценивается нормой Фробениуса разности матриц:
    \[
    \mathcal{D}_{PF}(Y^{(1)}, Y^{(2)}) = ||\mathbf{T}^{(1)} - \mathbf{T}^{(2)}||_F =
     \sqrt{\sum_{i,j=1}^{K} (T_{ij}^{(1)} - T_{ij}^{(2)})^2}.
    \]
\end{enumerate}
\subsection*{Второй этап отбора: расстояние Шамфера}
Для фильтрации выбросов, прошедших первый этап отбора, используется \textbf{расстояние Шамфера}
 \cite{Wu2021}, которое измеряет геометрическое расхождение между аттракторами
  в фазовом пространстве. В отличие от эргодической метрики, оно чувствительно к форме и
   плотности распределения траекторий, позволяя отфильтровать системы с похожими
    операторами, но отличной геометрией аттрактора.

Расстояние Шамфера вычисляется между \textbf{реконструированными аттракторами} временных рядов.
 Пусть $\{x^{(1)}_t\}_{t=1}^{N_1}$ и $\{x^{(2)}_t\}_{t=1}^{N_2}$ — исходные одномерные временные ряды.
  Согласно теореме Такенса \cite{Takens1981}, их динамика может быть восстановлена в пространстве
   вложения $\mathbb{R}^m$ путем построения векторов состояния:
\begin{align}
\mathbf{y}^{(1)}_i &= \bigl[x^{(1)}_i, x^{(1)}_{i+\tau}, x^{(1)}_{i+2\tau}, \dots, x^{(1)}_{i+(m-1)\tau}\bigr], \quad i = 1, \dots, N_1-(m-1)\tau, \\
\mathbf{y}^{(2)}_j &= \bigl[x^{(2)}_j, x^{(2)}_{j+\tau}, x^{(2)}_{j+2\tau}, \dots, x^{(2)}_{j+(m-1)\tau}\bigr], \quad j = 1, \dots, N_2-(m-1)\tau,
\end{align}
где $m$ — размерность вложения, $\tau$ — временная задержка. В данной работе использовались
 параметры реконструкции $m=3$ и $\tau=10$. Выбор размерности $m=3$ соответствует истинной размерности фазового пространства системы Лоренца, а
 значение задержки $\tau=10$ является распространенным выбором для данной системы Совокупности этих векторов образуют
 множества точек аттракторов:
\begin{equation}
A_1 = \{\mathbf{y}^{(1)}_i\} \subset \mathbb{R}^m, \qquad A_2 = \{\mathbf{y}^{(2)}_j\} \subset \mathbb{R}^m.
\end{equation}

Расстояние Шамфера между аттракторами определяется как симметричная метрика, усредняющая 
минимальные евклидовы расстояния от каждой точки одного облака к ближайшей точке другого:
\begin{equation}
\mathcal{D}_{\text{Chamfer}}(A_1, A_2) = \underbrace{\frac{1}{|A_1|}\sum_{\mathbf{x}\in A_1}\min_{\mathbf{y}\in A_2}\|\mathbf{x}-\mathbf{y}\|_2}_{\text{направление } A_1 \to A_2} + \underbrace{\frac{1}{|A_2|}\sum_{\mathbf{y}\in A_2}\min_{\mathbf{x}\in A_1}\|\mathbf{y}-\mathbf{x}\|_2}_{\text{направление } A_2 \to A_1},
\end{equation}
где $|A_k|$ обозначает мощность множества $A_k$.

Двунаправленная природа метрики гарантирует чувствительность к форме,
 локальной плотности и топологической структуре аттракторов, позволяя
 эффективно различать динамические режимы (например, хаос и предельные циклы), даже при
  формальной близости параметров системы.

\textbf{Двухуровневый отбор} работает следующим образом: сначала по метрике $\mathcal{D}_{PF}$ 
отбирается топ-$N_1$ кандидатов, затем для них вычисляется $\mathcal{D}_{\text{Chamfer}}$, и 
финальный отбор производится из топ-$N_2$ рядов с наименьшим геометрическим расхождением.

\paragraph{Процедура эксперимента.}
Процедура эксперимента включала следующие шаги. Сначала было сгенерировано 10,000
 вспомогательных рядов системы Лоренца, где управляющий параметр $r$ случайным 
 образом варьировался в диапазоне $r \in [23, 33]$ вокруг эталонного значения $r_0=28$. 
 На первом этапе отбора каждый из этих рядов-кандидатов был сравнен с эталонным рядом с 
 помощью метрики $D_{PF}$, в результате чего были отобраны топ 480 рядов с наименьшим 
 значением нормы Фробениуса. На втором этапе для этой группы было вычислено расстояние $D_{\text{Chamfer}}$ и выбраны 
топ 48 рядов с минимальным геометрическим расхождением, которые составили группу Отобранная аугментация.

Для контроля была сформирована вторая группа, Случайная аугментация, в которую вошли 48 рядов,
 выбранных случайным образом из исходного пула кандидатов. Выбор размера групп был обусловлен количеством
  ядер, доступных для параллельной обработки на суперкомпьютере.

Далее для каждого ряда из обеих групп была сформирована гибридная обучающая выборка, состоящая из 10000 
точек оригинального ряда ($r=28$) и 20,000 точек вспомогательного. Качество прогноза для каждой такой выборки
 оценивалось по метрикам RMSE, MAPE и доле непрогнозируемых точек. Для наглядности эксперимент проводился 
 с параметром $\alpha=10$, который контролирует строгость отбора прогнозируемых точек. Это привело к увеличению
  их доли и, как следствие, к снижению RMSE, что позволило более четко визуализировать различия. Контрольные 
  расчеты с исходным параметром $\alpha=10/3$ показали аналогичные по сути выводы.
\subsubsection{Анализ результатов}

\begin{figure}[htbp]
    \centering
    \includegraphics[width=\textwidth]{Augmentations_benefit_visualisation.png} % Замените путь!
    \caption{Сравнение метрик качества прогноза для выборок, аугментированных по
     двухуровневому критерию (верхний ряд) и случайными рядами (нижний ряд).
      Горизонтальные линии показывают эталонные уровни для выборок из $N_{base}
      $ и $N_{base}+N_{aux}$ (точки) наблюдений оригинального ряда. Видно,
       что двухуровневый отбор существенно снижает разброс метрик и исключает катастрофические выбросы.}
    \label{fig:two_level_validation}
\end{figure}

Результаты эксперимента, представленные на Рисунке~\ref{fig:two_level_validation},
 демонстрируют высокую эффективность двухуровневого критерия отбора. Визуально заметно,
  что разброс значений метрик для отобранной группы значительно меньше, а сами значения
   в среднем лучше, чем для случайной группы. Ключевым достижением является
    \textbf{уменьшение количества серьезных выбросов ошибки} (RMSE > $0.04$),
     наблюдаемых при использовании только метрики Перрона-Фробениуса.

Это визуальное наблюдение находит строгое количественное подтверждение в результатах 
статистического анализа. Для сравнения распределений метрик между группами \textbf{Отобранная аугментация} 
и \textbf{Случайная аугментация} был применен критерий Манна-Уитни.

Результаты статистического теста показали ($p$-value $Исходя из требования\approx 0$) 
для всех ключевых метрик: RMSE, MAPE и доли непрогнозируемых точек (которая для отобранной
 группы снизилась в среднем на $\approx 60\%$ относительно базовой выборки).
Крайне низкие $p$-values ($p \approx 0$) однозначно свидетельствуют о том, что
 двухуровневый подход неслучайно выделяет ряды-доноры, улучшающие качество прогноза.

\textbf{Влияние второго уровня (расстояние Шамфера)} особенно заметно при сравнении 
с результатами одноуровневого отбора. Включение геометрической метрики позволило отфильтровать ряды с
  похожими операторами, но непригодной для метода мотивов геометрией аттрактора.
  Это привело к снижению максимальной ошибки по RMSE с 0.076 до 0.056 и уменьшению дисперсии
   ошибки в 3.3 раза. При этом средняя и медианная точность (0.028 и 0.026 соответственно) с
   охранились на уровне, практически идентичном идеальному случаю обучения на выборке из 30,000 точек оригинального ряда (0.027).

Таким образом, результаты эксперимента подтверждают, что двухуровневый критерий,
 основанный на сочетании эргодической и геометрической метрик, является 
 эффективным инструментом отбора рядов-доноров. Он не только снижает число непрогнозируемых 
 точек, но и \textbf{стабилизирует качество прогноза}, исключая катастрофические выбросы.
  В дальнейшем можно рассмотреть варианты с увеличением размерности матрицы перехода при наличии
   более длинных рядов, что позволит повысить точность первого уровня отбора.